\section{What is resolved, what remains, and how to use the package}
\label{sec:status}

\subsection{The mathematical results and their exact reach}

\begin{longtable}{P{.24\linewidth}P{.39\linewidth}P{.28\linewidth}}
\toprule
Problem & Result in this report & Remaining condition\\
\midrule\endhead
Compatible exact approximations & A finite uniqueness theorem gives
coherence of an explicit PR producer. & Correctness and uniqueness must
be proved at each level.\\[4pt]
Implicit algebraic construction & Filtered correction solves a simultaneous
polynomial system in each nilpotent truncation with a fixed step count.
& A seed root and inverse Jacobian at the first level.\\[4pt]
Formal power-series solutions & Coherent exact jets, with constant seed and
uniqueness. & Analytic radii and tail bounds are separate estimates.\\[4pt]
Local modular roots & Coherent solutions modulo \(q^n\), including composite
\(q\), from a unit Jacobian seed. & No inference to global integer or
rational solvability.\\[4pt]
Differential equations over integral rings & Hurwitz coefficients are
constructed by a triangular primitive recurrence without factorial
division. & Analytic interpretation needs additional assumptions.\\[4pt]
Geometric zero-count input & Sharp norm and subspace fiber bounds for
\(Y^m-cX^d\) with lower-weight terms and coprime \(m,d\). & The geometric
curve class is explicit; a general divisor library is not supplied.\\[4pt]
Finite spaces of functions & An explicit monomial basis and exact eventual
dimension \(L+1-(m-1)(d-1)/2\). & No identification with arbitrary
Riemann--Roch spaces is asserted.\\[4pt]
Finite-cover argument & A cover by \(M\) proper subspaces yields a bound
\(ML\) on a supplied distinct point list. & A proof or construction of the
cover is still needed.\\[4pt]
Deep arithmetic formalization & Compiled core and object composition;
explicit new targets and dependency order. & Full embedded lifting and
norm derivations remain unfinished.\\
\bottomrule
\end{longtable}

The norm certificate is produced from the coefficient tables. The jet and
modular root producers are constructed from their finite seeds. Thus the
applications do more than name a proof field containing the conclusion.
On the other hand, the finite-cover corollary is deliberately conditional:
its covering hypothesis is substantial mathematics, and a finite checker
for a proposed cover does not prove that an adequate cover always exists.

\subsection{Three boundaries that cannot be erased by a representation}

\paragraph{A count bound is not a cutoff.}
For a machine code \(e\), let \(R(e,t)\) say that the machine first halts
at exactly step \(t\). This is a PR predicate, and determinism proves
that at most one \(t\) satisfies it. If there were a total computable
bound \(B(e)\) containing every such witness, checking the first
\(B(e)\) steps would decide halting. Hence no uniform conversion from
an arbitrary finite-cardinality bound to a search cutoff exists, even
for PR families of sets with at most one element. This is a general
obstruction to an inference, not an independence result for any specific
Diophantine theorem. A new height theorem for a particular curve family
could still supply the missing information.

\paragraph{Finite Cauchy approximations do not decide exact algebraic cases.}
Define a real \(a_e\) to be \(2^{-t}\) if machine \(e\) first halts
at step \(t\), and zero if it never halts. At precision \(k\), simulate
only \(k\) steps and return the corresponding dyadic value if halting
has occurred, otherwise zero. This is a uniform PR rational evaluator;
its error is at most \(2^{-k}\). Externally, in the ordinary classical metatheory, every \(a_e\) is rational, yet a
uniform exact zero test for these names would decide halting. A claim
that every Cauchy-presented rational can be uniformly converted to an
exact rational code with decidable equality is therefore false.

Our regular lifting theorem uses a different input: an exact seed and
a finite invertibility certificate. For \(Y^2-X=0\) at \((0,0)\),
the derivative with respect to \(Y\) is zero, so the hypothesis fails.
The theorem does not silently select a singular Puiseux branch.

\paragraph{Coherent local roots need not be global points.}
The polynomial \(Z^2-2\) has the root \(3\) modulo \(7\), and its
derivative there is \(6\), a unit. The modular lifting theorem constructs
compatible roots modulo every \(7^n\). It has no rational root: writing
a putative root in coprime integers leads to both numerator and denominator
being even. Thus the tower supplies local arithmetic information, not an
integer or rational witness. Likewise a formal identity is not already
a quantified analytic error bound, even when an independent convergence
theorem is available for the same data.

\subsection{Consequences for the unresolved repository requests}

The new norm theorem settles the zero-count component of the Siegel
finite-cover route for the stated weighted curves. The dimension and
function-field constructions reduce further algebraic setup in that class.
The quantitative approximation theorem, suitable covers, local estimates,
and the global height/list problem remain. This does not prove the full
Siegel or Faltings theorem in \ipra.

The lifting theorem provides a common finite backend for the regular
parts of formal implicit construction, modular arithmetic and local
expansions. It does not construct all singular Newton--Puiseux branches
from arbitrary Cauchy coefficients. The existing exact-algebraic Puiseux
argument and analytic regular-chart estimates remain relevant separate
modules.

For Mordell--Weil the obstacle identified in the inspected packet is not
ordinary height descent. Its supplied representative list must still be
constructed and proved to cover the desired quotient. The checked descent
core in this archive can consume such data, but is not a resolution of
that missing arithmetic existence theorem. For Chebotarev,
Siegel--Walfisz and Bombieri--Vinogradov the analytic estimates and their
uniform witnesses similarly remain outside the results proved here.

\subsection{Reproducibility and the next concrete formalization steps}

The archive includes the complete LaTeX source and compiled PDF,
standalone Lean 4.32.0 sources, exact-arithmetic example programs, source
provenance and machine-readable validation records. The programs check
finite instances and invariants. Their passing results are regression
evidence, not proofs of the universal algebraic claims.

The first useful extension is to finish the common integer, rational,
list, polynomial and finite determinant library inside the object theory.
The next is the generic bounded-loop compiler and its coefficient-prefix
invariants. With those in place, frozen-inverse lifting has a short proof
dependency graph: polynomial Taylor remainder, Jacobian congruence,
residual improvement, uniqueness and restriction. The weighted norm
development then reuses the same polynomial kernel and adds a finite
determinant library, rescaling, Vandermonde factors and finite-list
grouping. Each endpoint is a \texttt{Prf PRAAx} object, transported to the
native calculus only afterwards.

The point of the framework is that these are common reusable modules.
The full logic of finite observation, the data consumed by a construction,
and the exact theorem remaining to be proved are visible at each boundary.
That makes future successes composable and prevents a failed existence
argument from disappearing inside a convenient representation.
